% Generated draft: requires mathematical and source review.
\documentclass[11pt]{article}
\usepackage[margin=25mm]{geometry}
\usepackage{fontspec}
\setmainfont{Latin Modern Roman}
\usepackage{amsmath,amssymb,mathtools,unicode-math}
\setmathfont{Latin Modern Math}
\usepackage{enumitem,needspace,xurl,hyperref}
\setlength{\parindent}{0pt}
\setlength{\parskip}{5pt}
\setlength{\emergencystretch}{3em}
\newcommand{\N}{\mathbb{N}}
\newcommand{\Z}{\mathbb{Z}}
\newcommand{\Q}{\mathbb{Q}}
\newcommand{\R}{\mathbb{R}}
\newcommand{\C}{\mathbb{C}}
\newcommand{\F}{\mathbb{F}}
\newcommand{\A}{\mathbb{A}}
\newcommand{\PP}{\mathbb{P}}
\newcommand{\E}{\mathbb{E}}
\newcommand{\eps}{\varepsilon}
\newcommand{\dd}{\,\mathrm{d}}
\begin{document}
\section*{UnsolvedMath ID: 30003164}
\textbf{Source problem code:} catalogue entry
\subsection{Definitions and mathematical statement}
Let \(k\) be a perfect field of characteristic \(p>0\), let \(X\) be a smooth \(d\)-dimensional variety over \(k\), and let \(Y\hookrightarrow X\) be a normal-crossing divisor. Write \(\mathcal I_Y\subseteq\mathcal O_X\) for its ideal sheaf. The usual infinitesimal thickenings are \(Y_s=V(\mathcal I_Y^s)\), indexed by integers \(s\geq1\); this is an explicit interpretation of the source's unspecified thickening convention.

Write \(\mathcal K^M_{n,X\mid Y_s}\) for the relative Milnor \(K\)-theory coefficient sheaf denoted by \(K_n^M(X,Y_s)\) in the supplied record. The source does not specify its construction. In particular, it does not say whether this means a kernel of restriction of Milnor \(K\)-sheaves or a relative sheaf defined by a modulus condition, nor does it specify the version of Milnor \(K\)-theory used on the possibly nonreduced schemes \(Y_s\). These choices must be fixed before the question is a fully determined mathematical assertion.

For \(\tau\in\{\mathrm{Zar},\mathrm{Nis}\}\), let \(H^i_\tau(X,-)\) denote degree-\(i\) sheaf cohomology on the small Zariski or Nisnevich site of \(X\), respectively. For compatible choices of the relative coefficient sheaves on these sites, change of topology gives the comparison maps in the following question.

\textbf{Comparison question, with unresolved coefficient and index conventions.} Is the canonical map
\[
\varprojlim_{s\geq1} H^i_{\mathrm{Zar}}\!\left(X,\mathcal K^M_{n,X\mid Y_s}\right)
\longrightarrow
\varprojlim_{s\geq1} H^i_{\mathrm{Nis}}\!\left(X,\mathcal K^M_{n,X\mid Y_s}\right)
\]
an isomorphism for every \(k,X,Y\) satisfying the stated hypotheses and all intended indices \(i,n\)? Here the displayed limits are interpreted as ordinary inverse limits, with transition maps induced by the inclusions \(Y_s\hookrightarrow Y_{s+1}\). The record does not explicitly specify the ranges of \(i,n\), or whether its limit notation was intended to denote ordinary inverse limits or comparison of inverse systems as pro-objects.

The source singles out the case in which \(k\) is finite and \(i=n=d\), attributing the conjectured comparison to Kato and Saito. It does not clarify whether that attribution covers the entire displayed generality or specifically this distinguished case. No properness or projectivity assumption on \(X\) is stated in the supplied record.
\subsection{Short English statement}
The question asks whether relative Milnor \(K\)-theory cohomology, taken over increasingly thick neighborhoods of a normal-crossing divisor, has the same inverse limit in the Zariski and Nisnevich topologies. It particularly highlights degree \(d\) and Milnor \(K\)-degree \(d\) over a finite field. The coefficient sheaf, index ranges, and intended meaning of the limit require clarification. The supplied open status is a dated curation claim, not an independently verified conclusion.
\subsection{Sources}
\begin{itemize}
\item Supplied local UnsolvedMath registry record, permanent ID \(30003164\), source problem code \texttt{catalogue entry}, title: Zariski--Nisnevich Comparison for Relative Milnor \(K\)-Theory. The record belongs to the Oberwolfach Reports Open Problems collection and attributes the conjecture to Kato and Saito; no individual proposer is supplied.
\item Supplied report citation: \emph{Algebraic K-theory and Motivic Cohomology} (2017), Oberwolfach Reports, DOI 10.4171/owr/2016/31, pp.\ 1753--1807. \url{https://doi.org/10.4171/owr/2016/31}. The registry separately records proposed year 2016; the relationship between these year labels has not been independently checked.
\item Supplied literature source: Oberwolfach report on Kato--Saito theory. \url{https://ems.press/content/serial-article-files/46635}.
\item Supplied literature source: \emph{K-theory and logarithmic Hodge--Witt theory}. \url{https://arxiv.org/abs/1512.04703}.
\item Literature assessment dated exactly \texttt{2026-08-21}: the supplied curation labels the problem open and states that the full Kato--Saito comparison remains conjectural in the stated generality, while reciprocity-sheaf and \(p\)-adic results establish important special cases without proving the extracted assertion. This is supplied dated curation, not an independently verified literature assessment. The background expressly requests independent verification of time-sensitive status and cited claims.
\end{itemize}
\end{document}
